\documentclass[reqno]{amsart}
\usepackage[T1]{fontenc}
\usepackage{amsmath,amssymb,amsthm}
\usepackage{xr-hyper}
\usepackage[hidelinks]{hyperref}
\externaldocument[dp-]{20261126T030900--denominators-of-projections-to-lifts__workbook}
\externaldocument[tw-]{20261126T031000--twists-with-few-congruences__workbook}
\externaldocument[w2-]{20261126T060600--lengths-as-values-of-polynomials__workbook}

\newcounter{chap}
% Chapter numbers, generated from 20261125T012000--grh-workbook__guide.org
% by scripts/make-chapters.py.  Do not edit: change the order of the
% index and rerun the script (./preview-site.sh and CI do this).
%
% Usage in a sheet, after \newcounter{chap}:
%   % Chapter numbers, generated from 20261125T012000--grh-workbook__guide.org
% by scripts/make-chapters.py.  Do not edit: change the order of the
% index and rerun the script (./preview-site.sh and CI do this).
%
% Usage in a sheet, after \newcounter{chap}:
%   % Chapter numbers, generated from 20261125T012000--grh-workbook__guide.org
% by scripts/make-chapters.py.  Do not edit: change the order of the
% index and rerun the script (./preview-site.sh and CI do this).
%
% Usage in a sheet, after \newcounter{chap}:
%   \input{chapters}\chapnumber{<sheet ID>}
\makeatletter
\providecommand\chapters@set[2]{\@namedef{chapters@#1}{#2}}
\chapters@set{20261125T001000}{1}% zeros-to-ramp-sums
\chapters@set{20261125T001100}{2}% inert-primes-and-estimate-b
\chapters@set{20261125T005000}{3}% galois-fourier-and-averaging
\chapters@set{20261125T005500}{4}% columns-from-a-sphere
\chapters@set{20261125T000000}{5}% slope-method
\chapters@set{20261126T030100}{6}% a-twist-for-each-prime
\chapters@set{20261125T003000}{7}% the-chain-and-the-square-root
\chapters@set{20261125T010000}{8}% rank-lemma-and-zero-estimate
\chapters@set{20261125T010500}{9}% poincare-series-and-near-orthogonality
\chapters@set{20261124T223000}{10}% elliptic-curves-for-the-proof
\chapters@set{20261125T009500}{11}% fock-model-and-theta-intertwiners
\chapters@set{20261124T223100}{12}% complex-multiplication-and-inert-reduction
\chapters@set{20261125T008000}{13}% one-cm-determinant
\chapters@set{20261124T214500}{14}% supersingular-theta-divisibility
\chapters@set{20261124T214600}{15}% the-isotropic-plane
\chapters@set{20261125T013700}{16}% pairings-on-a-formal-group-of-height-two
\chapters@set{20261125T013500}{17}% hecke-nondegeneracy-on-the-tree
\chapters@set{20261126T030200}{18}% the-bound-is-attained
\chapters@set{20261125T009000}{19}% kronecker-limit-and-chowla-selberg
\chapters@set{20261125T002000}{20}% definite-theta-lifts
\chapters@set{20261125T015500}{21}% siegel-weil-in-the-smallest-case
\chapters@set{20261125T002100}{22}% definite-theta-lifts-and-norms
\chapters@set{20261125T016500}{23}% hecke-operators-and-theta-lifts
\chapters@set{20261125T016700}{24}% divisibility-along-the-chain
\chapters@set{20261125T004000}{25}% denominators-and-sha
\chapters@set{20261125T004100}{26}% tate-twists-and-sha
\chapters@set{20261126T020430}{27}% the-minus-part-by-gauss-sums
\chapters@set{20261125T014100}{28}% congruence-modules
\chapters@set{20261125T014500}{29}% eisenstein-congruences-and-selmer-classes
\chapters@set{20261125T014700}{30}% pseudocharacters-and-the-length-bound
\chapters@set{20261125T011000}{31}% galois-representations-black-box
\chapters@set{20261126T030300}{32}% the-data
\chapters@set{20261126T020210}{33}% bessel-poincare-series-and-their-gram-operators
\chapters@set{20261126T030400}{34}% the-seesaw-identity
\chapters@set{20261126T030500}{35}% divisibility-of-the-composite-lift
\chapters@set{20261126T030600}{36}% theorem-a-short-forms
\chapters@set{20261126T030700}{37}% the-determinant-of-a-chain
\chapters@set{20261126T030800}{38}% the-cost-of-one-lift
\chapters@set{20261126T030900}{39}% denominators-of-projections-to-lifts
\chapters@set{20261126T031000}{40}% twists-with-few-congruences
\chapters@set{20261126T031100}{41}% theorem-c-and-the-main-inequality
\providecommand\chapnumber[1]{%
  \@ifundefined{chapters@#1}%
    {\PackageError{chapters}{Sheet #1 is not listed in the index}%
      {List it in 20261125T012000--grh-workbook__guide.org\MessageBreak
       and run scripts/make-chapters.py.}}%
    {\edef\thechap{\@nameuse{chapters@#1}}}}
\makeatother
\chapnumber{<sheet ID>}
\makeatletter
\providecommand\chapters@set[2]{\@namedef{chapters@#1}{#2}}
\chapters@set{20261125T001000}{1}% zeros-to-ramp-sums
\chapters@set{20261125T001100}{2}% inert-primes-and-estimate-b
\chapters@set{20261125T005000}{3}% galois-fourier-and-averaging
\chapters@set{20261125T005500}{4}% columns-from-a-sphere
\chapters@set{20261125T000000}{5}% slope-method
\chapters@set{20261126T030100}{6}% a-twist-for-each-prime
\chapters@set{20261125T003000}{7}% the-chain-and-the-square-root
\chapters@set{20261125T010000}{8}% rank-lemma-and-zero-estimate
\chapters@set{20261125T010500}{9}% poincare-series-and-near-orthogonality
\chapters@set{20261124T223000}{10}% elliptic-curves-for-the-proof
\chapters@set{20261125T009500}{11}% fock-model-and-theta-intertwiners
\chapters@set{20261124T223100}{12}% complex-multiplication-and-inert-reduction
\chapters@set{20261125T008000}{13}% one-cm-determinant
\chapters@set{20261124T214500}{14}% supersingular-theta-divisibility
\chapters@set{20261124T214600}{15}% the-isotropic-plane
\chapters@set{20261125T013700}{16}% pairings-on-a-formal-group-of-height-two
\chapters@set{20261125T013500}{17}% hecke-nondegeneracy-on-the-tree
\chapters@set{20261126T030200}{18}% the-bound-is-attained
\chapters@set{20261125T009000}{19}% kronecker-limit-and-chowla-selberg
\chapters@set{20261125T002000}{20}% definite-theta-lifts
\chapters@set{20261125T015500}{21}% siegel-weil-in-the-smallest-case
\chapters@set{20261125T002100}{22}% definite-theta-lifts-and-norms
\chapters@set{20261125T016500}{23}% hecke-operators-and-theta-lifts
\chapters@set{20261125T016700}{24}% divisibility-along-the-chain
\chapters@set{20261125T004000}{25}% denominators-and-sha
\chapters@set{20261125T004100}{26}% tate-twists-and-sha
\chapters@set{20261126T020430}{27}% the-minus-part-by-gauss-sums
\chapters@set{20261125T014100}{28}% congruence-modules
\chapters@set{20261125T014500}{29}% eisenstein-congruences-and-selmer-classes
\chapters@set{20261125T014700}{30}% pseudocharacters-and-the-length-bound
\chapters@set{20261125T011000}{31}% galois-representations-black-box
\chapters@set{20261126T030300}{32}% the-data
\chapters@set{20261126T020210}{33}% bessel-poincare-series-and-their-gram-operators
\chapters@set{20261126T030400}{34}% the-seesaw-identity
\chapters@set{20261126T030500}{35}% divisibility-of-the-composite-lift
\chapters@set{20261126T030600}{36}% theorem-a-short-forms
\chapters@set{20261126T030700}{37}% the-determinant-of-a-chain
\chapters@set{20261126T030800}{38}% the-cost-of-one-lift
\chapters@set{20261126T030900}{39}% denominators-of-projections-to-lifts
\chapters@set{20261126T031000}{40}% twists-with-few-congruences
\chapters@set{20261126T031100}{41}% theorem-c-and-the-main-inequality
\providecommand\chapnumber[1]{%
  \@ifundefined{chapters@#1}%
    {\PackageError{chapters}{Sheet #1 is not listed in the index}%
      {List it in 20261125T012000--grh-workbook__guide.org\MessageBreak
       and run scripts/make-chapters.py.}}%
    {\edef\thechap{\@nameuse{chapters@#1}}}}
\makeatother
\chapnumber{<sheet ID>}
\makeatletter
\providecommand\chapters@set[2]{\@namedef{chapters@#1}{#2}}
\chapters@set{20261125T001000}{1}% zeros-to-ramp-sums
\chapters@set{20261125T001100}{2}% inert-primes-and-estimate-b
\chapters@set{20261125T005000}{3}% galois-fourier-and-averaging
\chapters@set{20261125T005500}{4}% columns-from-a-sphere
\chapters@set{20261125T000000}{5}% slope-method
\chapters@set{20261126T030100}{6}% a-twist-for-each-prime
\chapters@set{20261125T003000}{7}% the-chain-and-the-square-root
\chapters@set{20261125T010000}{8}% rank-lemma-and-zero-estimate
\chapters@set{20261125T010500}{9}% poincare-series-and-near-orthogonality
\chapters@set{20261124T223000}{10}% elliptic-curves-for-the-proof
\chapters@set{20261125T009500}{11}% fock-model-and-theta-intertwiners
\chapters@set{20261124T223100}{12}% complex-multiplication-and-inert-reduction
\chapters@set{20261125T008000}{13}% one-cm-determinant
\chapters@set{20261124T214500}{14}% supersingular-theta-divisibility
\chapters@set{20261124T214600}{15}% the-isotropic-plane
\chapters@set{20261125T013700}{16}% pairings-on-a-formal-group-of-height-two
\chapters@set{20261125T013500}{17}% hecke-nondegeneracy-on-the-tree
\chapters@set{20261126T030200}{18}% the-bound-is-attained
\chapters@set{20261125T009000}{19}% kronecker-limit-and-chowla-selberg
\chapters@set{20261125T002000}{20}% definite-theta-lifts
\chapters@set{20261125T015500}{21}% siegel-weil-in-the-smallest-case
\chapters@set{20261125T002100}{22}% definite-theta-lifts-and-norms
\chapters@set{20261125T016500}{23}% hecke-operators-and-theta-lifts
\chapters@set{20261125T016700}{24}% divisibility-along-the-chain
\chapters@set{20261125T004000}{25}% denominators-and-sha
\chapters@set{20261125T004100}{26}% tate-twists-and-sha
\chapters@set{20261126T020430}{27}% the-minus-part-by-gauss-sums
\chapters@set{20261125T014100}{28}% congruence-modules
\chapters@set{20261125T014500}{29}% eisenstein-congruences-and-selmer-classes
\chapters@set{20261125T014700}{30}% pseudocharacters-and-the-length-bound
\chapters@set{20261125T011000}{31}% galois-representations-black-box
\chapters@set{20261126T030300}{32}% the-data
\chapters@set{20261126T020210}{33}% bessel-poincare-series-and-their-gram-operators
\chapters@set{20261126T030400}{34}% the-seesaw-identity
\chapters@set{20261126T030500}{35}% divisibility-of-the-composite-lift
\chapters@set{20261126T030600}{36}% theorem-a-short-forms
\chapters@set{20261126T030700}{37}% the-determinant-of-a-chain
\chapters@set{20261126T030800}{38}% the-cost-of-one-lift
\chapters@set{20261126T030900}{39}% denominators-of-projections-to-lifts
\chapters@set{20261126T031000}{40}% twists-with-few-congruences
\chapters@set{20261126T031100}{41}% theorem-c-and-the-main-inequality
\providecommand\chapnumber[1]{%
  \@ifundefined{chapters@#1}%
    {\PackageError{chapters}{Sheet #1 is not listed in the index}%
      {List it in 20261125T012000--grh-workbook__guide.org\MessageBreak
       and run scripts/make-chapters.py.}}%
    {\edef\thechap{\@nameuse{chapters@#1}}}}
\makeatother

\ifcsname chapters@20261126T060700\endcsname
  \chapnumber{20261126T060700}
\else
  \renewcommand{\thechap}{?}
\fi
\renewcommand{\thesection}{\thechap.\arabic{section}}
\theoremstyle{definition}
\newtheorem{problem}{Problem}[section]
\newtheorem{definition}{Definition}[section]
\theoremstyle{plain}
\newtheorem{theorem}{Theorem}[section]
\newcommand{\Z}{\mathbb{Z}}
\newcommand{\Q}{\mathbb{Q}}
\newcommand{\F}{\mathbb{F}}
\newcommand{\Oc}{\mathcal{O}}
\newcommand{\ord}{\mathrm{ord}}
\newcommand{\rk}{\operatorname{rk}}
\newcommand{\len}{\operatorname{len}}
\newcommand{\Fr}{\mathrm{Fr}}
\newcommand{\sla}{\mathrm{sl}}
\newcommand{\ev}{\mathrm{ev}}
\newcommand{\Rel}{\mathrm{Rel}}
\newcommand{\Mir}{\mathrm{Mir}}
\newcommand{\tot}{\mathrm{tot}}
\newcommand{\cyc}{\chi_{\mathrm{cyc}}}
\newcommand{\Kc}{\mathcal{K}}
\newcommand{\Xc}{\mathcal{X}}
\setlength{\emergencystretch}{3em}

\title{Tied primes}
\author{}
\date{November 26, 2026}

\begin{document}
\maketitle
\refstepcounter{chap}\label{chap}

\paragraph{Prerequisites.}
From Chapter~\ref{w2-chap}: Definitions~\ref{w2-def:chi}, \ref{w2-def:selmer-datum} and~\ref{w2-def:tied}, Problem~\ref{w2-pr:normalize}, and \(g^{+} = g\), \(g^{-} = (g^c)^{-1}\).  Section~\ref{dp-sec:statement}: \(G_E\), \(g^c\), \(\cyc\), tame generators, coefficient rings.  Cocycles, inertia and Frobenius \cite[Ch.~I, IV, VII]{SLF}; Kummer theory \cite[Ch.~IV, \S3]{Ne}; ideles and the Artin map \cite[Ch.~VI]{Ne}; lengths.

\medskip
The twisting theorem of Chapter~\ref{tw-chap} uses the polynomial input of Chapter~\ref{w2-chap}, which takes as given the bounded tied primes (Section~\ref{w2-sec:data}): boundedly many auxiliary primes, with ties, bound each tied Selmer length by \(\mathrm{glob} + \mathrm{plc}\).  Why is the number of primes bounded independently of the level \(\ell^{h_1}\) of their characters?  One minimizes the sum of the tied lengths over data and levels, then compares a tied module with a module of relations.

Throughout, \(E\), \(\Sigma_0\), \(g_0\) are as in Section~\ref{w2-sec:universal}; \(T_{\mathrm{Sel},1}, \dots, T_{\mathrm{Sel},n}\) are Selmer data for \((\Sigma_0, g_0)\) over a coefficient ring \(\Oc\) with uniformizer \(\pi\), all using the same fixed place above each finite place of \(E\).  In this sheet they are written \(T_1, \dots, T_n\), and \(T\) stands for any one of them; the target forms do not occur.  \(\gamma = \gamma(T)\) and \(\lambda\) are the inertia element and scalar of a datum, and \(\len\) is \(\Oc\)-length.  For \(0 \ne a \in \Oc\), \(\bar T_a = T/aT\) as in Definition~\ref{w2-def:selmer-datum} and \(\bar T_a^{*} := \mathrm{Hom}_{\Z}(\bar T_a, \Q/\Z)\), with \(\Oc\) acting by precomposition.  We fix integers \(c \ge 0\), \(h_1 > c\) and a neighbourhood \(W\) of \(g_0\); \(\pi^{ce_{\Oc}}\Oc = \ell^c\Oc\).  For tied primes \(Q\) (Definition~\ref{w2-def:tied}) and \(0 \ne a \mid \ell^c\), \(Z_Q(T_i, a)\) is the tied module of \(T_i\) at level \(a\).  Commutators are \([g, h] := ghg^{-1}h^{-1}\).

\section{A prime with a character}\label{sec:prime}

In the examples of this section \(E = \Q(i)\), \(\ell = 3\).  A character for \(\{w\}\) to \(\Z/3^{h_1}\) with value \(1\) at \(s_w\) needs \(3^{h_1} \mid \#k_w - 1\) and a surjective logarithm on \(k_w^\times\) killing the residues of \(i\), \(3\).

\begin{problem}\label{pr:sixty-one}
\textup{(a)} For \(h_1 = 1\): the least prime \(p \equiv 1 \pmod{12}\) with \(3^{(p-1)/3} \equiv 1 \pmod p\) is \(61 = (6+5i)(6-5i)\); modulo \(w = (6+5i)\), \(2\) is a primitive root, \(i \equiv 11 = 2^{15}\), \(3 \equiv 2^{6}\), and both are cubes.  \textup{(b)} For \(h_1 = 2\): the least prime \(p \equiv 1 \pmod{36}\) with \(3^{(p-1)/9} \equiv 1 \pmod p\) is \(1621\); modulo \(w = (39 - 10i)\), \(i \equiv 166 = 2^{405}\), \(3 \equiv 2^{1548}\), and both are ninth powers.
\end{problem}

\begin{problem}\label{pr:kummer}
Let \(w\) be a finite place of a number field \(E\) with \(\#k_w = 1 + \ell^{h_1}k'\), and \(T'\) a finite set of finite places, \(w \notin T'\), such that every finite idele is a principal idele times an idele that is a unit at every place outside \(T'\).  Suppose every element of \(E^\times\) that is a unit outside \(T'\) is an \(\ell^{h_1}\)-th power modulo \(w\), and let \(\log \colon k_w^\times \to \Z/\ell^{h_1}\) be a surjective homomorphism.  There is exactly one homomorphism \(\omega\) from the finite ideles to \(\Z/\ell^{h_1}\), trivial on \(E^\times\), with \(\omega(x) = \log(x_w \bmod w)\) for every idele \(x\) that is a unit at each place outside \(T'\); it is continuous.
\end{problem}

Let \(E\) be CM, \(w \ne \bar w\) of prime norm, \(w \nmid \ell\), and \(T'\) contain the places above \(\ell\).  Then \(\lg_w(\mathrm{Art}\,x) := \omega(x/\bar x)\) defines a character for \(\{w\}\), by reciprocity \cite[Ch.~VI]{Ne}; normalize \(\log\) so that \(\lg_w(s_w) = 1\).  With the Artin map of Definition~\ref{dp-def:aux-twist}, at a principal place \(v = (\pi_v) \ne w, \bar w\) its value at an arithmetic Frobenius element is \(\log(\pi_v/\bar\pi_v \bmod w)\).

\begin{problem}\label{pr:frobenius-value}
Keep \(\ell = 3\), \(h_1 = 1\), \(w = (6+5i)\) and the primitive root \(2\).  \(3^{24} \equiv 1 \pmod{73}\), \(73 = (8+3i)(8-3i)\); modulo \(w\), \((8+3i)/(8-3i) \equiv 2^{40}\) and \((8-3i)/(8+3i) \equiv 2^{20}\).  The values of \(\lg_w\) at arithmetic Frobenius elements at the two primes above \(73\) are \(1\) and \(-1\), in some order.
\end{problem}

\emph{Chebotarev's theorem, existence form} (black box (i)).  Let \(W' \subseteq G_E\) be nonempty and a neighbourhood of each of its points, and \(S\) a finite set of finite places.  There are a place \(v \notin S\) of prime norm and a place \(P\) of \(\bar{\Q}\) above \(v\) such that some arithmetic Frobenius element at \(P\) lies in \(W'\), and every arithmetic Frobenius element at \(P\) lies in \(W'\).  \emph{Status}: classical \cite[Ch.~VII, Thm.~13.4]{Ne}.

\begin{problem}\label{pr:correct}
Let \(\Fr\) be an arithmetic Frobenius element at a place \(P\) above \(w\), \(s_w\) its tame generator, \(\lg_w\) a homomorphism to \(\Z/\ell^{h_1}\) with \(\lg_w(s_w) = 1\), and \(e' \in \Z\) with \(e' \bmod \ell^{h_1} = \lg_w(\Fr)\).  Then \(\Fr\, s_w^{-e'}\) is an arithmetic Frobenius element at \(P\) with \(\lg_w(\Fr\, s_w^{-e'}) = 0\).
\end{problem}

We use the following positive-level form of the construction of auxiliary primes in \cite{GL}.  \emph{Auxiliary primes.}  Let \(S\), \(R\) be finite sets of finite places, \(g\) an element of the normal subgroup generated by the inertia groups above the places of \(R\) outside \(S \cup \bar S\) and not above \(\ell\), and \(\Xc\) a neighbourhood of \(g\).  For every integer \(h_1 \ge 1\) there are a place \(v \notin R \cup \bar R \cup S \cup \bar S\), \(v \nmid \ell\), of prime norm with \(\bar v \ne v\) and \(\#k_v \equiv 1 \pmod{\ell^{h_1}}\); a place above \(v\) with its tame generator \(s_v\); an arithmetic Frobenius element \(\Fr_v \in \Xc\) there; and a locally constant homomorphism \(\lg_v \colon G_E \to \Z/\ell^{h_1}\) with \(\lg_v(s_v) = 1\), \(\lg_v(\Fr_v) = 0\), coming by \(x \mapsto x/\bar x\) from a character of the finite ideles trivial on \(E^\times\) and on the ideles that are units outside \(S \cup \bar S\) and the places above \(\ell\) and congruent to \(1\) at \(v\).

The auxiliary-prime statement gives a \emph{first prime}: for every \(h_1 > c\), tied primes \(Q = \{w_*\}\) for \(c\), \(W\), with all slopes \(0\).

\section{Adding a prime}\label{sec:adding}

\begin{definition}\label{def:candidate}
Let \(Q\) be tied primes at level \(\ell^{h_1}\), for \(c\) and \(W\), with first prime \(w_*\).  We define a \emph{candidate for row \(i\)} to consist of a place \(v\) and a choice of data at \(v\).  The norm of \(v\) is prime, \(v\ne\bar v\), \(v,\bar v\notin\Sigma_0\cup Q\), \(v\) does not lie above the rational prime \(\#k_{w_*}\), and \(\#k_v\equiv1\pmod{\ell^{h_1}}\).  The data consist of a place above \(v\), its tame generator \(s_v\), an arithmetic Frobenius element \(\Fr_v\) there, and a locally constant homomorphism \(\lg_v\colon G_E\to\Z/\ell^{h_1}\) with \(\lg_v(s_v)=1\) and \(\lg_v(\Fr_v)=0\).  For each \(D\subseteq(Q\smallsetminus\{w_*\})\cup\{u:u\mid\ell\}\), there must be a character \(\theta_D\) of the finite ideles to \(\Z/\ell^{h_1}\), trivial on \(E^\times\) and on the ideles that are units outside \(D\cup\bar D\) and congruent to \(1\) at \(v\), such that \(\lg_v(\mathrm{Art}\,x)=\theta_D(x/\bar x)\) for every finite idele \(x\).  Finally, \(\lg_w(\Fr_v)\) is \(1\) for \(w=w_*\) and \(0\) for the other \(w\in Q\), and \(T_i(\Fr_v^{\pm})\equiv1\pmod{\ell^c}\).  We denote a candidate with its data by \(v\).  For \(0 \ne a \mid \ell^c\), \(Z^v_Q(T_j, a)\) is defined as \(Z_Q(T_j, a)\) but with ramification allowed also at \(v\) and \(\bar v\), and no tie there.
\end{definition}

\emph{Tame local cocycles} (black box (iii)).  Let \(v \nmid \ell\) be a place at which \(T\) is unramified, \(P\) a place of \(\bar{\Q}\) above \(v\), and \(s \in I_P\) an element whose image topologically generates the tame quotient of the inertia group \(I_P\).  A locally constant cocycle \(z \colon G_E \to \bar T_a\) with \(z(s) = 0\) vanishes on every inertia group above \(v\); if moreover \(z(F) = 0\) for an element \(F\) of the decomposition group \(G_P\) mapping to a topological generator of \(G_P/I_P\), then \(z = 0\) on \(G_P\).  \emph{Status}: classical; the tame quotient is procyclic \cite[Ch.~IV, \S2]{SLF}.  If \(T\) is unramified also at \(\bar v\), the criterion there uses \(s_v^{-}\) and \(\Fr_v^{-}\); both hypotheses hold for a candidate.

\begin{definition}\label{def:evaluations}
For a candidate \(v\): \(\ev_s(z) := (z(s_v^{+}), z(s_v^{-}))\) and \(\ev_{\Fr}(z) := (z(\Fr_v^{+}), z(\Fr_v^{-}))\), in \(\bar T_a^{\,2}\).
\end{definition}

\begin{problem}\label{pr:old}
Let \(v\) be a candidate for row \(i\).  \textup{(a)} \(Q \cup \{v\}\), with any \(\sla_i(v) \in \Oc\) and \(\sla_j(v) = 0\) for \(j \ne i\), is tied primes with first prime \(w_*\).  \textup{(b)} On \(Z^v_Q(T_j, a)\), \(\ker \ev_s = Z_Q(T_j, a)\).  \textup{(c)} The tied module of \(T_i\) for \(Q \cup \{v\}\) is \(Z^v_Q(T_i, a) \cap \ker(\ev_s - \sla_i(v)\ev_{\Fr})\), and that of \(T_j\), \(j \ne i\), is \(Z_Q(T_j, a)\).  \textup{(d)} If all slopes of row \(j\) are \(0\), every member of \(Z_Q(T_j, a)\) is unramified outside \(\Sigma(T_j)\).
\end{problem}

\begin{problem}\label{pr:tilt-example}
On \(C = \F_5^2\) let \(\ev_s(u_1, u_2) = (u_1, 0)\) and \(\ev_{\Fr}(u_1, u_2) = (2u_1, u_2)\).  Then \(\#\ker(\alpha\,\ev_s - \ev_{\Fr}) = 1, 1, 5, 1, 1\) for \(\alpha = 0, \dots, 4\); the old module \(\ker \ev_s\) has \(5\) elements, and \(\ker \ev_s \cap \ev_{\Fr}^{-1}(\ev_s C) = \{0\}\).
\end{problem}

\begin{problem}\label{pr:tilting}
Let \(C\) be an \(\Oc\)-module of finite length and \(\ev_s, \ev_{\Fr} \colon C \to (\Oc/\pi^k)^{n'}\) linear.  Outside at most \(n'\) residue classes of \(\alpha \in \Oc\) modulo \(\pi\),
\[
  \len \ker(\alpha\,\ev_s - \ev_{\Fr}) \le \len\bigl(\ker \ev_s \cap \ev_{\Fr}^{-1}(\ev_s C)\bigr).
\]
\end{problem}

Locally constant cocycles of \(G_E\) in a finite module, unramified outside a fixed finite set, are finitely many \cite[Thm.~8.3.20]{NSW}; so these modules have finite length.  Define \(\tot(Q) := \sum_{j=1}^n \sum_{k=1}^{ce_{\Oc}} \len Z_Q(T_j, \pi^k)\), and \(\tot_\varnothing\) as the same sum for the modules of locally constant cocycles unramified outside \(\Sigma(T_j)\), with the local relations and \(z(\gamma(T_j)) = 0\).

\begin{problem}\label{pr:decrease}
Assume \(2\rk T_i \cdot ce_{\Oc} + 1 < \#k_{\Oc}\).  Let \(v\) be a candidate for row \(i\) with \(\ev_{\Fr}(z) \notin \ev_s Z^v_Q(T_i, \pi^k)\) for some \(1 \le k \le ce_{\Oc}\) and \(z \in Z_Q(T_i, \pi^k)\).  For some unit \(\alpha\), the tied primes \(Q \cup \{v\}\) with \(\sla_i(v) = \alpha^{-1}\) and \(\sla_j(v) = 0\) for \(j \ne i\) have \(\tot(Q \cup \{v\}) < \tot(Q)\).
\end{problem}

\begin{problem}\label{pr:minimum}
Assume the inequality of Problem~\ref{pr:decrease} for every \(i\), and fix a first prime \(w_*\).  Among all tied primes (sets, data and slopes) with first prime \(w_*\) and \(\#Q + \tot(Q) \le 1 + \tot(\{w_*\})\), take one with \(\tot(Q)\) least.  Then \(\#Q \le 1 + \tot_\varnothing\), independent of \(h_1\), \(w_*\); and \(\ev_{\Fr} Z_Q(T_i, \pi^k) \subseteq \ev_s Z^v_Q(T_i, \pi^k)\) for every \(i\), every \(1 \le k \le ce_{\Oc}\) and every candidate \(v\) for row \(i\).
\end{problem}

\section{The count}\label{sec:count}

\begin{problem}\label{pr:primitive}
Let \(k \le ce_{\Oc}\) be least with \(\pi^k Z_Q(T, \ell^c) = 0\).  Multiplication by \(\pi^{ce_{\Oc} - k}\) induces an isomorphism \(Z_Q(T, \pi^k) \simeq Z_Q(T, \ell^c)\), and if \(k \ge 1\), some value of some member of \(Z_Q(T, \pi^k)\) has a unit coordinate.
\end{problem}

\emph{The setting of the count.}  Let \(V\) be a finite discrete \(\Oc\)-module killed by \(\ell^{h_1}\) with a continuous \(\Oc\)-linear action of \(G_E\), \(V^{*} := \mathrm{Hom}_{\Z}(V, \Q/\Z)\), \(G'\) a normal subgroup of finite index, closed, fixing \(V\) and \(\mu_{\ell^{h_1}}\), and \(\Sigma'\) a finite set of finite places outside which \(V\) is unramified.  At each \(v \in \Sigma'\) fix a place of \(\bar{\Q}\) and an \(\Oc\)-module \(A_v\) of continuous cocycles on its decomposition group containing the coboundaries.  For a set \(R\) of places, \(Z(R)\) is the module of locally constant cocycles \(G_E \to V\), unramified outside \(\Sigma' \cup R\), with restriction in \(A_v\) at each \(v \in \Sigma'\); \(f_E\) is the number of infinite places of \(E\).

\begin{definition}\label{def:relations}
For a finite set \(\Sigma''\) of finite places, \(\Rel(\Sigma'')\) is the \(\Oc\)-module of locally constant \(y \colon G_E \to V^{*}\), zero off \(G'\), additive on \(G'\), zero on the inertia groups above the places outside \(\Sigma''\), with \(y(gug^{-1}) = \cyc(g)\,y(u) \circ g^{-1}\) for \(g \in G_E\), \(u \in G'\) (\(\cyc\) read modulo \(\ell^{h_1}\)), and such that \(\sum_j y(\Fr_j)(z(s_j)) = 0\) for every finite family of distinct places \(v_j \notin \Sigma''\), places above them with arithmetic Frobenius elements \(\Fr_j \in G'\) and elements \(s_j\) of their inertia groups with \(s_j(r_j)/r_j = \zeta\) for some \(r_j\) with \(r_j^{\ell^{h_1}} \in E^\times\) of valuation \(1\) at \(v_j\), \(\zeta\) one fixed primitive \(\ell^{h_1}\)-th root of unity; and every \(z \in Z(\{v_j\})\).
\end{definition}

\emph{The Greenberg--Wiles count} (black box (ii)).  Let \(E\) be totally complex, \(\sigma \in G_E\) fix \(\mu_{\ell^\infty}\), with \([\sigma, g] \in G'\) for all \(g\) and no nonzero fixed vector in \(V\), and let \(p \ge 1\) be an integer prime to \(\ell\), divisible by no place of \(\Sigma'\), with the inertia groups above the primes dividing \(p\) contained in \(G'\).  There is a finite \(\Sigma'' \supseteq \Sigma'\) containing no place above a prime dividing \(p\) with
\[
  \#Z(\varnothing)\cdot (\#V)^{\#\Sigma' + f_E} \le \#\Rel(\Sigma'')\cdot \#V \cdot \prod_{v \in \Sigma'} \#A_v .
\]
\emph{Status}: classical in origin \cite[Thm.~2.19]{DDT}, \cite[Prop.~1.6]{Wi}; this form is read from the proof \cite{GL}.

\emph{The tied setting.}  For \(T = T_i\), tied primes \(Q\) and \(a = \pi^k\), \(1 \le k \le ce_{\Oc}\): \(V = \bar T_a\); \(G'\) is the kernel of the action on \(T/\ell^cT\) and on \(\mu_{\ell^{h_1}}\); \(\Sigma' := \Sigma(T) \cup Q' \cup \bar Q'\) with \(Q' := Q \smallsetminus \{w_*\}\); at \(v \in \Sigma(T)\) use the fixed place above \(v\), at \(w \in Q'\) that of \(s_w\), at \(\bar w\) its conjugate; \(A_v\) is the module of restrictions of the locally constant cocycles \(G_E \to \bar T_a\) with the local relations of \(T\) and the ties at \(Q\).  Read from the proof \cite{GL}: \(\sum_{v \mid \ell} \len A_v \le (\#\{v \mid \ell\} + f_E)\len \bar T_a\), and \(\len A_v \le \len \bar T_a\) plus the local term at \(v\) of Definition~\ref{w2-def:selmer-datum}, for \(v \in \Sigma(T)\) prime to \(\ell\).  Also read from the proof: \emph{let \(\Sigma''\) be finite with \(w_*, \bar w_* \notin \Sigma''\); if \(v\) is a candidate for row \(i\) with \(v, \bar v \notin \Sigma''\) and \(\Fr_v, \Fr_v^c \in G'\), then for \(y \in \Rel(\Sigma'')\), \(z' \in Z^v_Q(T_i, a)\) and \(a' \in \Oc\)}
\[
  y(\Fr_v)(a'z'(s_v)) + y(\Fr_v^c)(a'z'(s_v^{-})) = 0 .
\]

\begin{problem}\label{pr:coboundaries}
\textup{(a)} \(Z_Q(T, a) \subseteq Z(\varnothing) \cap \ker(z \mapsto z(\gamma))\).  \textup{(b)} The coboundaries lie in \(Z(\varnothing)\), and evaluation at \(\gamma\) maps them onto \(\bar T_a\).  \textup{(c)} \(\len Z(\varnothing) \ge \len \bar T_a + \len Z_Q(T, a)\).
\end{problem}

\begin{problem}\label{pr:count-example}
Let \(E = \Q(\sqrt{-5})\), \(L = \Q(i, \sqrt 5) = E(i)\), \(\ell = 3\), \(\Oc = \Z_3\), and let \(G_E\) act on \(V = \F_3\) through \(\mathrm{Gal}(L/E) = \{1, \delta\}\) by the sign.  \textup{(a)} The discriminants of \(E\), \(L\) are \(-20\), \(400\) and the class numbers \(2\), \(1\); \(L/E\) is unramified at every finite place, and the locally constant cocycles \(G_E \to V\) unramified at every finite place are exactly the \(3\) coboundaries.  \textup{(b)} There is \(\sigma \in G_E\) restricting to \(\delta\) and fixing \(\mu_{3^\infty}\); with \(G' = G_{L(\sqrt{-3})}\), \(p = 7\), \(h_1 = 1\), \(\Sigma' = \varnothing\), the count gives a finite \(\Sigma''\) with \(3 \cdot 3 \le \#\Rel(\Sigma'') \cdot 3\).
\end{problem}

\begin{problem}\label{pr:cancel}
In the tied setting, the count with \(\sigma = \gamma\) and \(p = \#k_{w_*}\) gives a finite \(\Sigma'' \supseteq \Sigma'\) containing no place above \(p\), in particular neither \(w_*\) nor \(\bar w_*\), with \(\len Z_Q(T, a) \le \len \Rel(\Sigma'') + \mathrm{plc}(T)\).
\end{problem}

\section{The mirror}\label{sec:mirror}

\begin{definition}\label{def:commutators}
\(\Kc = \Kc(T)\) is the subgroup of \(G_E\) generated by the commutators \([g, \gamma]\) and \([g, \gamma^c]\), \(g \in G_E\).
\end{definition}

\begin{problem}\label{pr:commutator}
Let \(z \colon G_E \to \bar T_a\) be a cocycle with \(z(\gamma) = 0\).  \textup{(a)} \(z([g, \gamma]) = (1 - \lambda)z(g)\).  \textup{(b)} \(\Kc\) is normal and stable under \(g \mapsto g^c\), acts trivially on \(T\), lies in \(G'\), and \(z(ghg^{-1}) = g\,z(h)\) for \(h \in \Kc\).  \textup{(c)} If some value of \(z\) has a unit coordinate, \(z(\Kc)\) spans \(\bar T_a\) over \(\Oc\).
\end{problem}

\begin{definition}\label{def:mirror}
For \(z \colon G_E \to \bar T_a\), \(y \colon G_E \to \bar T_a^{*}\) and \(a' \in \Oc\), define the \emph{mirror form} by \(\Mir_{a'}(h, h') := y(h')(a'z(h)) - y(h^c)(a'z(h'^c))\), with values in \(\Q/\Z\).
\end{definition}

From the statement on auxiliary primes the proof obtains, for tied primes \(Q\), a row \(i\), \(h \in \Kc(T_i)\), a finite set \(R\) and a neighbourhood \(\Xc\) of \(s_{w_*}h\): \emph{a candidate \(v\) for row \(i\) with \(v, \bar v \notin R\) and \(\Fr_v \in \Xc\)} (read from the proof \cite{GL}).

\begin{problem}\label{pr:diagonal}
Let \(z \in Z_Q(T_i, \pi^k)\) with \(\ev_{\Fr}(z) \in \ev_s Z^v_Q(T_i, \pi^k)\) for every candidate \(v\) for row \(i\), and \(y \in \Rel(\Sigma'')\) with \(w_*, \bar w_* \notin \Sigma''\).  Then \(\Mir_{a'}(h, h) = 0\) for all \(h \in \Kc(T_i)\) and \(a' \in \Oc\).
\end{problem}

\begin{problem}\label{pr:polarize}
Let \(z\) be a cocycle with \(z(\gamma) = 0\), \(y \in \Rel(\Sigma'')\), and suppose \(\Mir_{a'}(h, h) = 0\) for all \(h \in \Kc\) and \(a' \in \Oc\).  \textup{(a)} \(\Mir_{a'}(h, h') + \Mir_{a'}(h', h) = 0\) for \(h, h' \in \Kc\).  \textup{(b)} For \(h \in \Kc\), \(h' \in G_E\): \(\Mir_{a'}(\gamma h\gamma^{-1}, h') = \Mir_{\lambda a'}(h, h')\) and \(\Mir_{\lambda a'}(h', \gamma h\gamma^{-1}) = \Mir_{a'}(h', h)\).  \textup{(c)} \(\Mir_{a'} = 0\) on \(\Kc \times \Kc\) for every \(a'\).
\end{problem}

\begin{problem}\label{pr:mirror-example}
In the model group \(\F_5^2\) with \((u_1, u_2)^c = (u_2, u_1)\) and \(z(u) = u_1\), the linear \(y_0(u) = Au_1 + Bu_2\) with \(y_0(h)z(h) = y_0(h^c)z(h^c)\) for all \(h\) are exactly those with \(A = 0\); then \(y_0(h^c) = B\,z(h)\).
\end{problem}

\begin{problem}\label{pr:inject}
Fix a cocycle \(z \colon G_E \to \bar T_a\) such that \(z(\Kc)\) spans \(\bar T_a\), and assume that for every \(y \in \Rel(\Sigma'')\) the mirror form of \(z, y\) vanishes on \(\Kc \times \Kc\) for every \(a'\).  \textup{(a)} For each \(y\) there is a unique \(\Oc\)-linear \(\eta_y \colon \bar T_a \to \bar T_a^{*}\) with \(\eta_y(z(h)) = y(h^c)\) for \(h \in \Kc\).  \textup{(b)} \(\eta_y\) is \(G_E\)-equivariant into \(\bar T_a^{*c}\) of Definition~\ref{w2-def:selmer-datum}.  \textup{(c)} \(y \mapsto \eta_y\) is \(\Oc\)-linear and injective, and \(\len \Rel(\Sigma'') \le \mathrm{glob}(T)\).
\end{problem}

\begin{theorem}[tied primes]\label{thm:tied}
Let \(T_{\mathrm{Sel},1}, \dots, T_{\mathrm{Sel},n}\) be Selmer data for \((\Sigma_0, g_0)\) over \(\Oc\), \(c \ge 0\) an integer with \(2\rk T_{\mathrm{Sel},i}\,\ord_{\Oc}(\ell^{c}) + 1 < \#k_{\Oc}\) for all \(i\), and \(W\) a neighbourhood of \(g_0\) in \(G_E\).  There is a nonnegative integer \(n_Q\) such that for every integer \(h_1 > c\) there are tied primes at level \(\ell^{h_1}\), for \(c\) and \(W\), with \(\#Q \le n_Q\), for which the tied module of each \(T_{\mathrm{Sel},i}\) at level \(\ell^{c}\) has \(\Oc\)-length at most \(\mathrm{glob}(T_{\mathrm{Sel},i}) + \mathrm{plc}(T_{\mathrm{Sel},i})\).
\end{theorem}

Problems~\ref{pr:minimum}, \ref{pr:primitive}, \ref{pr:cancel}, \ref{pr:commutator}, \ref{pr:diagonal}, \ref{pr:polarize} and~\ref{pr:inject} prove Theorem~\ref{thm:tied}, with \(n_Q = 1 + \tot_\varnothing\).  Taken as given: (i), (ii), (iii), reciprocity, finiteness of cocycles; and, read from the proof, the auxiliary primes of Section~\ref{sec:prime}, the local bounds and the identity for a conjugate pair in the tied setting, and the candidates near \(s_{w_*}h\) of Section~\ref{sec:mirror}.  It gives the bounded tied primes of Section~\ref{w2-sec:data} when \(\Sigma_0\) contains the places above \(\ell\), as for the twisting theorem of Chapter~\ref{tw-chap}.

\paragraph{In the Lean formalization.}
\begin{sloppypar}
In the Lean proof \cite{GL} (\nolinkurl{grh/src/}; names beginning with \texttt{SelmerDatum.} are in \texttt{New.Selmer}), Theorem~\ref{thm:tied} is the case of \texttt{New.Selmer.exists\_isTiedPrimes\_card\_le\_and\_length\_le} (\nolinkurl{New/Congruences/Selmer/Descent.lean}) for the coefficient rings of Section~\ref{dp-sec:statement} and \(\Sigma_0\) containing the places above \(\ell\); Problem~\ref{pr:tilting} is the case of \texttt{New.Selmer.tilting}.  Definition~\ref{def:candidate} is the trivial-action-modulo-\(\ell^c\) case of \texttt{New.Selmer.IsNewAuxiliaryPrime}; \(\Kc\) is \texttt{SelmerDatum.sahGroup}; the mirror form is \texttt{SelmerDatum.mirrorForm}.  The statement on auxiliary primes is taken from the proof of \texttt{New.Selmer.exists\_isAuxPrimeDataAt}, and its use before Problem~\ref{pr:diagonal} is \texttt{SelmerDatum.exists\_isNewAuxiliaryPrime}.  (ii) applies \texttt{New.RelationSetting.exists\_card\_cocycles\_le} with \(m = \ell^{h_1}\); the tied setting, its local bounds and Problem~\ref{pr:cancel} follow the proof of \texttt{SelmerDatum.exists\_length\_tiedSelmer\_le}; the identity for a conjugate pair is \texttt{SelmerDatum.dualSelmer\_frob\_add\_conj\_frob\_eq\_zero}.  Other steps are in the proofs in \nolinkurl{Descent.lean}, \nolinkurl{Count.lean}, \nolinkurl{Mirror.lean}.
\end{sloppypar}

\section*{Hints}

\subsection*{Section~\ref{sec:prime}}
\begin{description}
\item[\ref{pr:sixty-one}] Check \(13\), \(37\).
\item[\ref{pr:frobenius-value}] \(i \equiv 11 \pmod w\).
\item[\ref{pr:correct}] \(s_w\) lies in the inertia group.  First prime: \(S = \varnothing\), \(g = g_0\), \(\Xc = W\), \(R \supseteq\) places over primes under \(\Sigma_0\).
\end{description}

\subsection*{Section~\ref{sec:adding}}
\begin{description}
\item[\ref{pr:old}] For (a), the character of \(v\) is trivial on \(E_w^\times\) for \(w \in Q \smallsetminus \{w_*\}\).  For (b), (c), (iii) at \(v\), \(\bar v\); for (d), at each \(w \in Q\) and \(\bar w\).
\item[\ref{pr:tilting}] Let \(C'\) be the sum of the kernels, \(\Lambda \subseteq \Oc^{n'}\) the preimage of \(\ev_s C'\); lift maps with source the free \(\Lambda\), and use a characteristic polynomial modulo \(\pi\).
\item[\ref{pr:decrease}] Problem~\ref{pr:tilting} at every level, \(n' = 2\rk T_i\); Problem~\ref{pr:old}; for \ref{pr:minimum}, \ref{pr:old}\,(d).
\end{description}

\subsection*{Section~\ref{sec:count}}
\begin{description}
\item[\ref{pr:primitive}] The \(\pi^k\)-torsion of \(\Oc/\ell^c\).
\item[\ref{pr:coboundaries}] The slope at \(w_*\) is \(0\); Problem~\ref{w2-pr:normalize}.
\item[\ref{pr:count-example}] For (a), an unramified cocycle restricts to a homomorphism of \(G_L\).  For (b), the quadratic subextension of \(E(\mu_{3^\infty})/E\) is \(E(\sqrt{-3})\).
\item[\ref{pr:cancel}] Check the hypotheses of (ii); at a tied place the tie and (iii) inject \(A_v\) into \(\bar T_a\) by the value at a Frobenius element.
\end{description}

\subsection*{Section~\ref{sec:mirror}}
\begin{description}
\item[\ref{pr:commutator}] For (c), reduce modulo \(\pi\); use (a) and that the \(T(g)\) span.
\item[\ref{pr:diagonal}] Compare with \(\Mir_{a'}(g, g)\), \(g = s_{w_*}h\); use a candidate outside \(\Sigma''\) near \(g\) and the conjugate-pair identity.
\item[\ref{pr:polarize}] Expand \(\Mir(hh', hh')\); for (c) use (a), (b); \(\lambda\), \(\lambda^2 - 1\) are units.
\item[\ref{pr:inject}] For (c), \(g = [g, \gamma]\cdot\gamma g\gamma^{-1}\).
\end{description}

{\small
\begin{thebibliography}{NSW}
\bibitem[DDT]{DDT} H.~Darmon, F.~Diamond and R.~Taylor, \emph{Fermat's Last Theorem}, Current Developments in Mathematics \textbf{1995} (1995), no.~1, 1--154.
\bibitem[GL]{GL} L.~Sernau, \emph{GRH for Dirichlet \(L\)-functions, in Lean~4}, a shorter rewriting of the Lean proof, text by Claude, \url{https://github.com/lsernau/grh-lean}, commit \texttt{72fead1}, readable edition \texttt{grh/}.
\bibitem[Ne]{Ne} J.~Neukirch, \emph{Algebraic Number Theory}, Grundlehren Math.\ Wiss.\ 322, Springer, 1999.
\bibitem[NSW]{NSW} J.~Neukirch, A.~Schmidt and K.~Wingberg, \emph{Cohomology of Number Fields}, 2nd ed., Grundlehren Math.\ Wiss.\ 323, Springer, 2008.
\bibitem[SLF]{SLF} J.-P.~Serre, \emph{Local Fields}, Grad.\ Texts in Math.\ 67, Springer, 1979.
\bibitem[Wi]{Wi} A.~Wiles, \emph{Modular elliptic curves and Fermat's Last Theorem}, Ann.\ of Math.\ \textbf{141} (1995), 443--551.
\end{thebibliography}
}

\end{document}
